\documentclass[12pt]{article}

\usepackage{../sty_files/packages}
\usepackage{../sty_files/shortcuts}
\usepackage{asymptote}


\usepackage{parskip}

\setlist[itemize]{label=$\bullet$}

% \newcommand{\Newdot}{{\leavevmode\put(0,.63){\circle*{2.5}}}}
%--------------------------------------------------------------%
\begin{document}

\begin{titlepage}
   \begin{center}
       \vspace*{1cm}
       \textbf{\Huge {No equality}}

       \vspace{1cm}
       \textbf{\Large{এস এম এ নাহিয়ান}}
       
       \vfill
       \includegraphics[width=0.4\textwidth]{images/meta.png}

       \large{গণিতযজ্ঞ}
   \end{center}
\end{titlepage}

\tableofcontents
%---------------------------------------------------------------%
\newpage

\section{Starting}
অসমতার ধারণা মূলত বাস্তব সংখ্যার অর্ডার বা ক্রম রয়েছে এই ফ্যাক্ট থেকে এসেছে। যেমন, $a > b$ হলে বাস্তবসংখ্যার
লাইনে $a$ সবসময় $b$ এর ডানে অবস্থান করবে। 
আমরা সমীকরণ নিয়ে কাজ করার সময় অনেক গুলো ব্যাসিক রুল মেনে ক্যালকুলেশন করি। তেমনি ইনইকুয়ালিটির 
ম্যানিপুলেশনের জন্যও কিছু ব্যাসিক রুল রয়েছে। আমি সাজেস্ট করবো বিগিনারদের এই রুল গুলো ভালোভাবে ফিল করার
ট্রাই করতে। রুলগুলো অনেক সিলি মনে হলেও জটিল ক্যাল্কুলেশনের মাঝে অসাবধান থাকলে যে কেউ এগুলো ভুল করে 
প্রবলেমকে ফেকসলভ করে ফেলতে পারে (আমি নিজেও অনেকবার করেছি)। 
\begin{theom*} %{Basic Rules for Inequalities}
    \begin{itemize}
        \item $x \geq y$ এবং $y \geq z$ হলে $x \geq z$ হবে। আর কোনোটিতে $>$ সম্পর্ক থাকলে $x > z$ হবে।
        \item $x \geq y$ হলে ধনাত্মক $k$ এর জন্য, $xk \geq yk$ এবং ঋনাত্মক $k$ এর জন্য $xk \leq yk$ হবে। আর, 
            $x > y$ হলে ধনাত্মক $k$ এর জন্য, $xk > yk$ এবং ঋনাত্মক $k$ এর জন্য $xk < yk$ হবে।
            \begin{description}
                \item[সতর্কতাঃ] অসমতার দুই পাশে কোনো ভ্যারিয়াবল গুণ করতে গেলে খেয়াল রাখতে হবে যেনো 
                সেটির রেঞ্জে ধনাত্মক আর ঋনাত্মক দুটোই যেনো একসাথে না থাকে।
            \end{description}
        \item $x \geq y$ এবং $a \geq b$ হলে $a + x \geq b + y$ হবে। আর কোনোটিতে $>$ সম্পর্ক থাকলে 
            $a + x > b+y$ হবে। খেয়াল করো, ২য় অসমতার যায়গায় সমীকরণ $a = b$ থাকলেও এভাবে যোগ করতে পারতে।
            \begin{description}
                \item[সতর্কতাঃ] অসমতার যোগ করা গেলেও তুমি এটা থেকে $x-a \geq y-b$ লিখতে পারবে না 
                কারণ এটি সবসময় সত্য হয় না। তবে কোনো অসমতা থেকে তার উলটো অসমতা বিয়োগ করা যায় 
                সেক্ষেত্রে প্রথম অসমতার চিহ্ন ব্যাবহার হবে। অর্থাৎ, $x \geq y,\ b \leq a \implies x-b \geq y-a$ 
                কিন্তু সেটি করা আর সোজা অসমতা দুটো যোগ করে পক্ষান্তর করা একই কাজ। 
            \end{description}
        \item $x \geq y$ এবং $x,y\neq 0$ হলে $x,\ y$ এর চিহ্ন একই হলে $\dfrac 1x \leq \dfrac 1y$ হবে। আর চিহ্ন 
            উলটো হলে $\dfrac 1x \geq \dfrac 1y$ হবে আর ইকুয়ালিটি থাকবে যদি এবং কেবল যদি $x=y$ হয়।
    \end{itemize}
\end{theom*}
\section{How to?}
\subsection{How to approach an inequality problem?}
সাধারণত, একটি ইনইকুয়ালিটি স্টেটমেন্ট $P_1$ প্রমাণের জন্য তুমি চেষ্টা করতে থাকবে সেটিকে সিমপ্লিফাই করে আরেকটি 
স্টেটমেন্ট $P_2$ (অবশই সেটিও সত্য হতে হবে) খুজে বের করতে যেনো, $P_2 \implies P_1$ হয় এরপর সেটি প্রমাণের চেষ্টা করবে। 
এভাবে $P_a$ থেকে $P_{a+1}$ এ যাওয়ার প্রোসেস চালাতে থাকবে যতোক্ষন না পর্যন্ত এমন কোনো স্টেটমেন্ট $P_k$ পাচ্ছো যেটি হয়, well-known অথবা, সহজে সরাসরি বীজগানিতিক 
ক্যালকুলেশন করে প্রমাণ করে ফেলতে পারো। অর্থাৎ, তোমার প্রমাণ করার চেষ্টাটি হবে অনেকগুলো স্টেটমেন্টের 
একটা চেইন যেখানে শুরুর স্টেটমেন্টটি সত্য (সম্ভবত কোনো থিওরেম বা well-known ইনইকুয়ালিটি) এবং প্রতিটি স্টেটমেন্ট তার পরের স্টেটমেন্টকে ইমপ্লাই করে আর 
সবার শেষ স্টেটমেন্টটি হবে সেটা যেটা তোমার প্রমাণ করা লাগবে। অনেকটা এরকম, 
$$P_k \implies P_{k-1} \implies \dots \implies P_2 \implies P_1$$
আর তুমি এই চেইনটি তৈরি করবে পেছন থেকে শুরু করে। 

তাহলে, কোনো ইনইকুয়ালিটি প্রমাণের জন্য তিনটি কাজ করতে পারো,
\begin{itemize}
    \item এমন কোনো অসমতা বের করা যেটি প্রমাণ করা আরো সহজ আর সেটি থেকে আমাদের প্রয়োজনীয় অসমতাটির পাওয়া যায়।
    \item প্রয়োজনীয় অসমতাটি কোনো well-known ইনইকুয়ালিটি তা উল্লেখ করে দেওয়া।
    \item সরাসরি বীজগানিতিক ক্যালকুলেশনের মাধ্যমে প্রমাণ কর। 
\end{itemize} 
তবে মনে রাখতে হবে, সবসময় এতো ডিরেক্ট ভাবে এপ্রোচ করা নাও যেতে পারে। কখনো একটা ইনইকুয়ালিটি প্রমানের জন্য সেটাকে 
দুইটা আলাদা ইনইকুয়ালিটি তে ভাগ করা সহজ হতে পারে। কখনো নতুন ভ্যারিয়াবল দিয়ে প্রতিস্থাপন করা লাগতে পারে। 
এছাড়াও আরো অনেক ক্লেভার আইডিয়াও কাজে লাগে।


% এখানে তৃতীয়টি নিয়ে বলার কিছু নেই। তবে প্রথম দুইটি এপ্লাই করার ক্ষেত্রে একটি অনেক কমন ট্রিক হলো, 
% পেছন থেকে চিন্তা করা। অর্থাৎ, আমরা শুরু করবো আমাদের যা প্রমান করতে হবে তা দিয়ে।
% এবং সেটি থেকে আমরা "যদি এবং কেবল যদি" সম্পর্কযুক্ত স্টেটমেন্টে যেতে থাকব যতোক্ষন না পর্যন্ত এমন কোনো 
% স্টেটমেন্টে যাবো যেটা প্রমান করা হয়তো আরো সহজ অথবা আমরা আগে থেকেই জানি সত্য। \\
% তবে এখানে খেয়াল রাখতে হবে প্রতিটা ধাপেই যেনো "যদি এবং কেবল যদি" সম্পর্ক বজায় থাকে। \\

\subsection{How to write a proof?}
কিছু বলার আগে চলো প্রমাণ লেখার দুটি উদাহরণ দেখে নিই,
\begin{exam} {উদাহরণ \thesection.\thesubsection.1}
    বাস্তব সংখ্যা $x$ এর জন্য প্রমাণ কর যে,
    $$x^2 + x > -1$$
\end{exam}
\begin{sol}
    \textbf{প্রমাণ ১ঃ} আমরা জানি, $\left(x+\frac 12\right)^2 \geq 0$। এখন,
    \begin{align*}
        \left(x+\frac 12\right)^2 &\geq 0\\
        \implies x^2 + x + \frac 14 &\geq 0\\
        \implies x^2 + x &\geq - \frac 14\\
    \end{align*}
    আর, $-\frac{1}{4} > -1$ সুতরাং, $x^2 + x > -1 \ \qed$ 
\end{sol}
\begin{sol}
    \textbf{প্রমাণ ২ঃ} এখানে, 
    \begin{align*}
        x^2 + x &> -1\\
        \iff x^2 + 2\cdot x\cdot \frac 12 &> -1\\
        \iff x^2 + 2\cdot x\cdot \frac 12 + \left(\frac 12\right)^2 &> -1 + \frac 14\\
        \iff \left(x+\frac 12\right)^2 &> -\frac{3}{4}\\
    \end{align*}
    যেটি সত্য কেননা বাস্তব সংখ্যার বর্গ সবসময় অঋণাত্মক হয়। সুতরাং আমাদের প্রথম স্টেটমেন্টটিও সত্য। $\qed$
\end{sol}
তো, একটা ইনইকুয়ালিটি প্রবলেম এর প্রমান তুমি সোজা ভাবেও লিখতে পারো আবার উলটো ভাবেও লিখতে পারো।
তবে কিছু বিষয়ে লক্ষ্য রাখতে হবে। যেমন, সোজা প্রমানটিতে আমি একটি well-known ইনইকুয়ালিটি থেকে শুরু 
করেছি আর প্রতি ধাপে $\implies$ সম্পর্কযুক্ত (যদিও এখানে সেগুলো "if and only if" সম্পর্কযুক্ত, তবে এই 
প্রমাণের জন্য শুধু $\implies$ সম্পর্কই যথেষ্ট) আরেকটি ইনইকুয়ালিটি তে গিয়েছি। আর উল্টোভাবে লিখতে গেলে 
প্রশ্নের অসমতাটি থেকে শুরু করে প্রতি ধাপে $\Longleftarrow$ অথবা, $\iff$ সম্পর্কযুক্ত 
অসমতায় গিয়েছি এবং শেষে একটি well-known ইনইকুয়ালিটি তে পৌছেছি। এখানে প্রতি স্টেপ এ $\iff$ সম্পর্ক 
রয়েছে তবে যদি শুধু $\Longleftarrow$ সম্পর্কযুক্ত 
হয় তাহলে সলুশন চেক করতে নিজের সমস্যা হবে।  
এখন প্রশ্ন করতে পারো উল্টোভাবে লেখা যদি সমস্যার সৃষ্টি 
করে তাহলে এভাবে লেখা কেনো শিখবো?\\
উত্তর হলো, পার্শিয়াল সলুশন লেখার জন্য। তুমি যখন একটা ইনইকুয়ালিটি প্রব্লেম পুরোপুরি সলভ করতে পারবে না তখন 
তুমি জানো না যে কি থেকে শুরু করলে সেই প্রবলেম স্টেটমেন্টটি পাওয়া যাবে। তখন তুমি পার্শিয়াল পয়েন্ট পাওয়ার জন্য 
প্রব্লেম স্টেটমেন্ট থেকে শুরু করে সিমপ্লিফাই করে অন্য কোনো সহজ ইনইকুয়ালিটি তে যাওয়ার চেষ্টা করতে থাকবে।
তাই সেক্ষেত্রে উল্টোভাবে প্রমাণ লেখাই বেটার হবে।


\newpage
\section{Some useful ineqs}
আমরা সবচেয়ে মৌলিক অসমতা টি দিয়েই শুরু করব। 
\subsection{The Trivial Inequality}
\subsubsection{Theorem Statement}
\begin{theom}{The Trivial Inequality}
    কোনো বাস্তব সংখ্যা $x$ এর জন্য, 
    $$x^2 \geq 0$$
    এবং সমতা ঘটে যদি এবং কেবল যদি $x = 0$ হয়।
\end{theom}
এটাকে আমরা একাধিক ভ্যারিয়েবলের জন্য এক্সটেন্ড করতে পারি,
\begin{theom}{The Trivial Inequality (Extended)}
    কয়েকটি বাস্তব সংখ্যা $x_1, \ x_2, \dots,\ x_n$ এর জন্য, 
    $$x_1^2 + x_2^2 + \dots + x_n^2 \geq 0$$
    এবং সমতা ঘটে যদি এবং কেবল যদি $x_1 = x_2 = \dots = x_n = 0$ হয়।
\end{theom}
এটা তোমরা সবাই অলরেডি জানো। কিন্তু তবুও উল্লেখ করার কারণ হলো, 
সাধাণত বেশিরভাগ অসমতা ই এই মৌলিক অসমতাটি থেকে আসে। তাই এটাকে 
"Mother of all inequality" ও বলা যায়।\\
\subsubsection{Examples}
এখন এটি ব্যাবহার করে কিছু উদাহরণ দেখা যাক, (এই উদাহরণ গুলো লেমা হিসেবে অনেক 
প্রব্লেম এ ব্যাবহার করা যায়) 
\begin{exam} {উদাহরণ \thesubsubsection.1}
    প্রমান করো যে, যেকোনো ধনাত্মক বাস্তব সংখ্যা $a$ এর জন্য, 
    $$a + \dfrac1a \geq  2$$
    এবং সমতা ঘটে যায় যদি এবং কেবল যদি $a = 1$ হয়। 
\end{exam}

\begin{sol}
    \textbf{সমাধানঃ} এখানে, $a > 0$ সুতরাং,
    \begin{align*}
        &a + \dfrac{1}{a} \geq 2\\
        \iff & a^2 + 1 \geq 2a\\
        \iff & a^2 - 2a + 1 \geq 0\\
        \iff & (a-1)^2 \geq 0
    \end{align*}
    এখানে প্রতি ধাপের সাথে পরের ধাপে যদি এবং কেবল যদি সম্পর্ক আছে। সুতরাং যেকোনো একটি স্টেটমেন্ট সত্য হলেই 
    সবগুলো সত্য হবে। আর Trivial Inequality তে $x = a-1$ বসিয়ে আমরা পাই, 
    শেষ স্টেটমেন্ট টি সত্য। সুতরাং, প্রথম স্টেটমেন্টটিও সত্য।
    আর শেষ স্টেটমেন্টে সমতা ঘটে, $a-1 = 0 \implies a = 1$ হলে। $\qed$
\end{sol}

\begin{exam} {উদাহরণ \thesubsubsection.2 (AM-GM Inequality for two varriables)}
    প্রমান করো যে, যেকোনো ধনাত্মক বাস্তব সংখ্যা $a, b$ এর জন্য, 
    $$\dfrac{a+b}{2} \geq \sqrt{ab}$$
    আর সমতা ঘটে যদি এবং কেবল যদি, $a = b$ হয়। 
\end{exam}
\begin{sol}
    \textbf{সমাধানঃ} Trivial inequality তে $x = (\sqrt a - \sqrt b)$ বসিয়ে পাই,
    \begin{align*}
        & (\sqrt a - \sqrt b)^2 \geq 0\\
        \iff & a - 2\sqrt{ab} +b \geq 0\\
        \iff &  a + b \geq 2\sqrt{ab}\\
        \iff & \dfrac{a+b}{2} \geq \sqrt{ab}
    \end{align*}
    আর সমতা ঘটে, $\sqrt a - \sqrt b = 0 \implies a = b$ হলে। $\qed$ 
\end{sol}


% \subsubsection{Excercise}

\subsection{AM-GM Inequality}
\subsubsection{Theorem Statement}
\begin{theom} {AM-GM Inequality}
    কতোগুলো ধনাত্মক বাস্তবসংখ্যা $x_1, \ x_2, \dots,\ x_n$ এর জন্য, 
    $$\dfrac{x_1+x_2+\dots + x_n}{n} \geq \sqrt[n]{x_1x_2\dots x_n}$$
    এবং সমতা ঘটে যদি এবং কেবল যদি, $x_1 = x_2 = \dots = x_n$ হয়।
\end{theom}
\begin{sol}
    \textbf{প্রমাণঃ} এটি আমরা প্রমান করবো কশি (Cauchy) ইন্ডাকশন দিয়ে। \\
    প্রথমেই বেস কেস, $n = 1$ এর জন্য ট্রিভিয়াল। আর $n = 2$ এর জন্য আমরা উদাহরণ ২.২ এ প্রমাণ করে এসেছি।\\
    এখন আমরা প্রমান করবো, কোনো $n = k$ এর জন্য এটি সত্য হলে $n = 2k$ এর জন্য ও সত্য। 
    আবার, কোনো $n = k$ এর জন্য সত্য হলে $n = k-1$ এর জন্য ও সত্য। তাহলেই AM-GM ইনইকুয়ালিটি প্রমান হয়ে যাবে
    সকল $n$ এর জন্য। (কেনো?) \\
    এখন ধরি, $n = k$ এর জন্য, $\dfrac{x_1+x_2+\dots + x_k}{k} \geq \sqrt[k]{x_1x_2\dots x_k}$। তাহলে,
    \begin{align*}
        \dfrac{x_1+x_2+\dots + x_{2k}}{2k} &= \dfrac{\dfrac{x_1+x_2+\dots + x_k}{k} + \dfrac{x_{k+1}+\dots + x_{2k}}{k}}{2}\\
        &\geq \dfrac{\sqrt[k]{x_1x_2\dots x_k} + \sqrt[k]{x_{k+1}\dots  x_{2k}}}{2}\\
        &\geq \sqrt{\sqrt[k]{x_1x_2\dots x_k} + \sqrt[k]{x_{k+1}\dots  x_{2k}}}\\
        &= \sqrt[2k]{x_1x_2\dots  x_{2k}}
    \end{align*}
    সুতরাং, $n = 2k$ এর জন্য ও এটি সত্য হবে।\\
    আবার, ধরি, $n = k$ এর জন্য, $\dfrac{x_1+x_2+\dots + x_k}{k} \geq \sqrt[k]{x_1x_2\dots x_k}$। তাহলে,
    এখন, $x_k = \dfrac{x_1+\dots + x_{k-1}}{k-1}$ বসিয়ে পাই,
    \begin{align*}
        \dfrac{x_1+x_2+\dots +x_{k-1} +\frac{x_1+\dots + x_{k-1}}{k-1}}{k} &\geq \sqrt[k]{x_1x_2\dots x_{k-1}\cdot \frac{x_1+\dots + x_{k-1}}{k-1}}\\
        \implies \dfrac{x_1+x_2+\dots +x_{k-1}}{k-1} &\geq \sqrt[k]{x_1x_2\dots x_{k-1}\cdot \frac{x_1+\dots + x_{k-1}}{k-1}}\\
        \implies \dfrac{x_1+x_2+\dots +x_{k-1}}{k-1} &\geq \left(x_1x_2\dots x_{k-1}\cdot \frac{x_1+\dots + x_{k-1}}{k-1}\right)^\frac{1}{k}\\
        \implies \dfrac{x_1+x_2+\dots +x_{k-1}}{k-1} &\geq \left(x_1x_2\dots x_{k-1} \right)^\frac{1}{k}  \left( \frac{x_1+\dots + x_{k-1}}{k-1}\right)^\frac{1}{k}\\
        \implies  \left( \frac{x_1+\dots + x_{k-1}}{k-1}\right)^\frac{k-1}{k} &\geq \left(x_1x_2\dots x_{k-1} \right)^\frac{1}{k} \\
        \implies  \left( \frac{x_1+\dots + x_{k-1}}{k-1}\right) &\geq \left(x_1x_2\dots x_{k-1} \right)^\frac{1}{k-1} \\
        \implies \dfrac{x_1+x_2+\dots + x_{k-1}}{k-1} &\geq \sqrt[k-1]{x_1x_2\dots x_{k-1}}\\
    \end{align*}
    সুতরাং ইন্ডাকশনের সাহায্যে, স্টেটমেন্টটি সকল $n \in \NN$ এর জন্যই সত্য। $\qed$
\end{sol}
\subsubsection{Examples}
\begin{exam} {উদাহরণ \thesubsubsection.1}
    ধনাত্মক বাস্তব সংখ্যা $a,b,c$ এর জন্য প্রমান কর যে,
    $$ a^2+b^2+c^2 \geq ab+bc+ca $$
\end{exam}
\begin{sol}
    \textbf{সমাধানঃ} এখানে আমরা সরাসরি $3$ ভ্যারিয়াবলের AM-GM ব্যাবহার করতে পারি না তবে $2$ ভ্যারিয়াবলের 
    AM-GM ব্যাবহার করে পাই, $$\dfrac{a^2 + b^2}{2} \geq ab$$ 
    একইভাবে, $$\dfrac{b^2 + c^2}{2} \geq bc \text{ এবং } \dfrac{c^2 + a^2}{2} \geq ca$$
    এখন, এই তিনটি ইনইকুয়ালিটিকে যোগ করে পেয়ে যাই,
    $$a^2+b^2+c^2 \geq ab+bc+ca \qed$$
\end{sol}
\begin{exam} {উদাহরণ \thesubsubsection.2}
    ধনাত্মক বাস্তব সংখ্যা $a,b,c$ এর জন্য প্রমান কর যে,
    $$ \dfrac{a^3}{bc} + \dfrac{b^3}{ca} + \dfrac{c^3}{ab} \geq a+b+c $$
\end{exam}
\begin{sol}
    \textbf{সমাধানঃ} এখানে, \begin{align*}
        \dfrac{a^3}{bc} + \dfrac{b^3}{ca} + \dfrac{c^3}{ab} &\geq a+b+c\\
        \iff \dfrac{a^4 + b^4 + c^4}{abc} &\geq a+b+c\\
        \iff a^4+b^4+c^4 &\geq a^2bc+b^2ca+c^2ab
    \end{align*}
    এখন, আমরা $4$ ভ্যারিয়াবলের AM-GM ব্যাবহার করে পাই, 
    $$\dfrac{a^4 + a^4 + b^4 + c^4}{4} \geq a^2bc$$ 
    একইভাবে, $$\dfrac{b^4 + b^4 + c^4 + a^4}{4} \geq b^2ca \text{ এবং } \dfrac{c^4 + c^4 + a^4 + b^4}{4} \geq c^2ab$$
    তিনটিকে যোগ করে পাই, 
    $$a^4+b^4+c^4 \geq a^2bc+b^2ca+c^2ab$$
    সুতরাং, প্রশ্নের স্টেটমেন্টটিও সত্য হবে। $\qed$ 
\end{sol}

\subsubsection{Excercise}
\begin{prob} {\thesubsubsection.1}
    বাস্তব সংখ্যা $a,b,c$ এর জন্য প্রমান কর যে,
    $$ a^3+b^3+c^3 \geq a^2b+b^2c+c^2a $$
\end{prob}

\begin{prob} {\thesubsubsection.2}
    $a,b,c$ ধনাত্মক বাস্তব সংখ্যা যেনো, $a^2 + b^2 + c^2 = 1$। তাহলে, $\dfrac{ab}{c} + \dfrac{bc}{a} + \dfrac{ca}{b}$ এর সর্বনিম্ন মান কত?
\end{prob}

% \subsection{Triangle Inequality}

\section{Example Problems}
\begin{exam} {উদাহরণ \thesection.1 (BdMO 2019 National HS P2)}
    প্রমাণ করো, $a,b,c$ ধনাত্মক বাস্তব সংখ্যা হলে,
    $$\dfrac{a}{bc} + \dfrac{b}{ca} + \dfrac{c}{ab} \geq \dfrac{2}{a} +\dfrac{2}{b}-\dfrac{2}{c}$$
\end{exam}
\begin{sol}
    \textbf{সমাধানঃ} এখানে, \begin{align*}
        \dfrac{a}{bc} + \dfrac{b}{ca} + \dfrac{c}{ab} &\geq \dfrac{2}{a} +\dfrac{2}{b}-\dfrac{2}{c}\\
        \iff \dfrac{a^2+b^2+c^2}{abc} &\geq \dfrac{2bc + 2ca - 2ab}{abc}\\
        \iff a^2 + b^2+c^2 &\geq 2bc+2ca - 2ab\\
        \iff a^2 + 2ab + b^2 + c^2 &\geq 2bc+2ca\\
        \iff(a+b)^2 + c^2 &\geq 2bc+2ca\\
        \iff(a+b)^2 - 2bc+2ca + c^2 &\geq 0\\
        \iff(a+b)^2 - 2c(a+b) + c^2 &\geq 0\\
        \iff(a+b-c)^2 &\geq 0\\
    \end{align*}
    যেটি অবশই সত্য। সুতরাং প্রথম অসমতাটিও সত্য। $\qed$
\end{sol}

\begin{exam} {উদাহরণ \thesection.2}
    প্রমান কর যে, $a,b,c$ তিনটি ধনাত্মক বাস্তব সংখ্যা হলে, $$6a^3 + 9b^3 +32c^3 + \dfrac 1{4abc}$$
    এর সর্বনিম্ন মান $6$। এবং $a,b,c$ এর কোন মানের জন্য সেটি পাওয়া যায়? 
\end{exam}
\begin{sol}
    \textbf{সমাধানঃ} ধরি, $x = \sqrt[3]{6}a$, $y = \sqrt[3]{9}b$ এবং $z = \sqrt[3]{32}c$। এখন, মানগুলো 
    প্রতিস্থাপন করে পাই, \begin{align*}
        6a^3 + 9b^3 +32c^3 + \dfrac 1{4abc} &\geq 6\\
        \iff x^3 + y^3 + z^3 + \dfrac{1}{4\cdot \frac{1}{\sqrt[3]{6}}x\cdot \frac{1}{\sqrt[3]{9}}y\cdot \frac{1}{\sqrt[3]{32}}z} &\geq 6\\
        \iff x^3 + y^3 + z^3 + \dfrac{1}{4\cdot \frac{1}{\sqrt[3]{6\times 9\times 32}}xyz} &\geq 6\\
        \iff x^3 + y^3 + z^3 + \dfrac{1}{4\cdot \frac{1}{12}xyz} &\geq 6\\
        \iff x^3 + y^3 + z^3 + \dfrac{3}{xyz} &\geq 6\\
    \end{align*}
    এখন, $x^3 + y^3 + z^3 + \dfrac{3}{xyz} \geq 3xyz + \dfrac{3}{xyz}$ যেখানে ইকুয়ালিটি থাকবে যদি এবং কেবল যদি $x=y=z$ হয়।\\
    আবার $3xyz + \dfrac{3}{xyz} = 3\left(xyz + \dfrac1{xyz}\right) \geq 3\times 2 = 6$ আর এখানে ইকুয়ালিটি হোল্ড করবে যদি 
    এবং কেবল যদি, $xyz = 1$ হয়।\\ 
    সুতরাং, অসমতাটি সত্য। আর ইকুয়ালিটি সত্য হবে যদি এবং কেবল যদি, $xyz = 1$ ও $x=y=z$ হয় অর্থাৎ 
    $x=y=z=1 \implies a=\frac 1{\sqrt[3]{6}}, b=\frac 1{\sqrt[3]{9}}, c=\frac 1{\sqrt[3]{32}}$ হলে। $\qed$
\end{sol}

\begin{exam} {উদাহরণ \thesection.3}
    ধনাত্মক সংখ্যা $x,y,z$ এর জন্য প্রমাণ করো, 
    $$\dfrac{2}{x+y} + \dfrac{2}{y+z} + \dfrac{2}{z+x} \geq \dfrac{9}{x+y+z}$$
\end{exam}
\begin{sol}
    \textbf{সমাধানঃ} ধরি, $a = x+y$, $b = y+z$ এবং $c = z+x$। এখন, 
    \begin{align*}
        \dfrac{2}{x+y} + \dfrac{2}{y+z} + \dfrac{2}{z+x} &\geq \dfrac{9}{x+y+z}\\
        \iff \dfrac{2}{a} + \dfrac{2}{b} + \dfrac{2}{c} &\geq \dfrac{9}{\frac{a+b+c}{2}}\\
        \iff \dfrac{2}{a} + \dfrac{2}{b} + \dfrac{2}{c} &\geq \dfrac{9\times 2}{a+b+c}\\
        \iff \dfrac{1}{a} + \dfrac{1}{b} + \dfrac{1}{c} &\geq \dfrac{9}{a+b+c}\\
        \iff \dfrac{ab+bc+ca}{abc} &\geq \dfrac{9}{a+b+c}\\
        \iff (ab+bc+ca) &\geq \dfrac{9abc}{a+b+c}\\
        \iff (ab+bc+ca)(a+b+c) &\geq 9abc\\
        \iff a^2b+abc+ca^2+ab^2+b^2c\qquad &\\
        +abc + abc+bc^2+c^2a &\geq 9abc\\
        \iff a^2b+b^2c+c^2a +ca^2+ab^2  +bc^2 &\geq 6abc\\
    \end{align*}
    এখন, AM-GM ব্যাবহার করে পাই, $a^2b+b^2c+c^2a \geq 3abc$ এবং $ca^2+ab^2+bc^2 \geq 3abc$। আর 
    যোগ করে পাই, $$a^2b+b^2c+c^2a +ca^2+ab^2  +bc^2 \geq 6abc$$
    সুতরাং, আমাদের প্রথম অসমতাটিও সত্য। $\qed$
\end{sol}

\section{Practice Probs}
\begin{prob} {\thesection.1}
    ধরো $a,b,c,d > 0$। প্রমাণ করো, 
    $$\dfrac 1a + \dfrac 1b + \dfrac 4c + \dfrac {16}d \geq \dfrac{64}{a+b+c+d}$$
\end{prob}

\begin{prob} {\thesection.2}
    ধরো তিনটি ধনাত্মক সংখ্যা $a,b,c$ যেনো, $abc = 1$। প্রমাণ করো,
    $$a^2+b^2+c^2 \geq a+b+c$$
\end{prob}

\begin{prob} {\thesection.3}
    ধরো, $ a_1, a_2,\ldots, a_n$ অঋণাত্মক বাস্তব সংখ্যা। প্রমাণ করো,
    $$\frac{1}{1+ a_1}+\frac{ a_1}{(1+ a_1)(1+ a_2)}+\frac{ a_1 a_2}{(1+ a_1)(1+ a_2)(1+ a_3)}+\cdots+\frac{ a_1 a_2\cdots a_{n-1}}{(1+ a_1)(1+ a_2)\cdots (1+ a_n)} \leq 1 $$
    এবং কখন সমান হয় তা বের করো।
\end{prob}

\end{document}